% established_factors.tex (Paper II)
% Structural prose for the PASS rows of audit_table.tex.  Each
% subsection takes one PASS row and explains what the script's output
% means structurally -- i.e. WHY the verified statement is the right
% statement, not just THAT it is verified.  The closing subsection
% (\autoref{subsec:established_necessity_vs_sufficiency}) collects the
% rows into the statement "established containment $\supseteq$"
% and names the gap to "equality $=$" that the open subproblems
% (\autoref{subsec:open_subproblems}) close.

\label{sec:established_factors}

The audit table (Table~\ref{tab:audit}) records the PASS rows
establishing sub-claims of Eq.~\eqref{eq:automorphism_conjecture},
together with one structural-consistency row (SU(3) representation
branch) that narrows the interpretation of the colour factor.  Each
is a sympy-verified statement on the bipartite octahedral lattice
$\mathcal{T}_\diamond^3$.  This section explains what these rows
mean structurally and how they combine into the containment
$\mathrm{Aut}(\mathcal{T}_\diamond^3, \mathcal{A}=1) \supseteq SO(3,1)
\times SU(3) \times SU(2) \times U(1)$ on the extended per-site
amplitude $\mathbb{C}^{12} = \mathbb{C}^2 \otimes \mathbb{C}^2
\otimes \mathbb{C}^3$.  The closing subsection
(\autoref{subsec:established_necessity_vs_sufficiency}) states the
gap between the verified containment and the conjectured equality;
that gap is the substance of Paper~II's three open subproblems
(\autoref{subsec:open_subproblems}).

\subsection{Discrete spatial automorphisms: $|\mathrm{Aut}_\text{discrete}(\Gamma, \mathbf{V})| = 48$}
\label{subsec:discrete_spatial}

The basis $\mathbf{V} = \{\pm\mathbf{V}_1, \pm\mathbf{V}_2,
\pm\mathbf{V}_3\}$ of the bipartite octahedral lattice has six
elements.  The linear maps $A \in GL(3,\mathbb{R})$ permuting
$\mathbf{V}$ as a set are parametrised by a permutation $\sigma \in
S_3$ of $(\mathbf{V}_1, \mathbf{V}_2, \mathbf{V}_3)$ together with
independent signs $s \in \{\pm 1\}^3$ on each basis vector.  This
gives $|S_3| \cdot |\mathbb{Z}_2^3| = 6 \cdot 8 = 48$ distinct
candidates, all of which preserve $\mathbf{V}$ as a set.  The script
\nolinkurl{automorphism_discrete.py} enumerates the 48 candidates,
verifies each preserves $\mathbf{V}$, checks closure under matrix
multiplication, and confirms that exactly 12 of the 48 elements
satisfy $A^\top A = I$ in the standard inner product on
$\mathbb{R}^3$.

The resulting group is the hyperoctahedral group
$B_3 \cong \mathbb{Z}_2 \wr S_3$, abstractly the cubic point group
$O_h$.  The realisation on $\mathbb{R}^3$ in the $\mathbf{V}$-basis
is not, however, the standard cubic-symmetry orthogonal
representation: only the 12-element orthogonal subgroup is metric-
preserving, while the remaining 36 elements are linear shears that
permute the basis labels at the cost of distorting angles.

Structurally, the full 48-element group is the relevant discrete
spatial symmetry for the framework: the bipartite tick rule treats
the RGB and CMY labels symmetrically and does not single out the
standard metric.  The 12-element orthogonal subgroup is the
metric-preserving part --- the genuine crystallographic point
symmetry of $\mathcal{T}_\diamond^3$ in the standard
$\mathbb{R}^3$ embedding --- and the standard metric is restored
in the continuum limit by the $O_h$-averaging argument of
Paper~I~§6 that produces $SO(3,1)$.  No piece of the 48 is
discarded; the non-orthogonal piece is the discrete record of the
fact that the lattice does not commit to a Euclidean metric at the
microscale.

\subsection{The RGB sublattice contributes only $\mathbb{Z}_3 \subset SU(3)$}
\label{subsec:rgb_z3_obstruction}

The discrete RGB symmetry $S_3$ of $(\mathbf{V}_1, \mathbf{V}_2,
\mathbf{V}_3)$ embeds in $U(3)$ as the six permutation matrices on
the kinetic-step weight space $(\alpha_1, \alpha_2, \alpha_3) \in
\mathbb{C}^3$.  Within $U(3)$, the determinant separates the six
permutations: the three even permutations have $\det = +1$ and lie
in $SU(3)$ as the alternating subgroup $A_3 \cong \mathbb{Z}_3$;
the three transpositions have $\det = -1$ and lie in $U(3) \setminus
SU(3)$.  The script \nolinkurl{automorphism_rgb_su3.py} verifies the
embedding and the determinant split explicitly.

This is the load-bearing structural observation that forces the
per-site $\mathbb{C}^3$ extension.  The conjecture's $SU(3)$ factor
has real dimension 8 and is non-abelian; $\mathbb{Z}_3$ is finite
and abelian; closure under multiplication does not extend a finite
abelian discrete group to a continuous non-abelian Lie group, and
the commutator $[A, B] = AB - BA$ on the three $\mathbb{Z}_3$
matrices either vanishes or lands back in $\mathbb{Z}_3$, so the
Lie algebra generated is itself abelian.  An $\mathfrak{su}(3)$
factor cannot be a Lie-theoretic continuation of the framework's
existing discrete colour symmetry.  It must come from an additional
per-site internal index that the existing $\mathbb{C}^2 = (\psi_R,
\psi_L)$ amplitude does not carry.

The natural candidate is a per-site $\mathbb{C}^3$ ``colour-memory''
amplitude $(c_1, c_2, c_3)$ where component $c_j$ records the
amplitude that the wavefunction's most recent RGB tick was along
$\mathbf{V}_j$.  The eight Gell-Mann generators
$\{\lambda_a/2\}_{a=1\ldots 8}$ acting on this $\mathbb{C}^3$
supply $\mathfrak{su}(3)$ by definition; the framework's
contribution is identifying this $\mathbb{C}^3$ as the
\emph{minimal} per-site internal extension the conjecture's
$SU(3)$ factor needs.  The existing $\mathbb{Z}_3$ remains as the
discrete subgroup of $SU(3)$ that permutes the colour basis ---
not refuted but recovered as a finite subgroup of the larger
continuous symmetry.

\subsection{$SO(3,1) \times U(1)$ on the existing per-site $\mathbb{C}^2$: dim 7}
\label{subsec:lorentz_u1_on_C2}

Four candidate Lie-algebra factors of the conjecture admit
generators on the existing per-site chirality amplitude
$\mathbb{C}^2 = (\psi_R, \psi_L)$: the Lorentz rotations
$J_a = \tfrac{1}{2}\sigma_a$, the Lorentz boosts $K_a =
\tfrac{i}{2}\sigma_a$, a proposed per-site $SU(2)$ with generators
$T_a = \tfrac{1}{2}\sigma_a$, and the per-site $U(1)$ phase
generator $Q = i I_2$.  The script
\nolinkurl{automorphism_direct_product.py} verifies that
$\mathfrak{so}(3,1)$ closes correctly on $(J_a, K_a)$ (the standard
Weyl-spinor relations $[J_a, J_b] = i\epsilon_{abc}J_c$,
$[K_a, K_b] = -i\epsilon_{abc}J_c$, $[J_a, K_b] = i\epsilon_{abc}
K_c$) and that $Q$ commutes with every generator.  So
$\mathfrak{so}(3,1) \oplus \mathfrak{u}(1)$ acts as a direct sum
on $\mathbb{C}^2$, with total real dimension $6 + 1 = 7$.

The script also verifies the literal matrix identity $T_a = J_a$
for every $a \in \{1, 2, 3\}$.  The proposed per-site $SU(2)$
generators on the chirality $\mathbb{C}^2$ are not just isomorphic
to the Lorentz rotation generators --- they are the same matrices.
As subgroups of $GL(2, \mathbb{C})$, the proposed per-site $SU(2)$
on this carrier is the Lorentz rotation subgroup of $SO(3,1)$.  A
direct-product decomposition $SO(3,1) \times SU(2)$ on the
existing $\mathbb{C}^2$ would double-count the rotations.

This is the second load-bearing structural observation, and it
parallels the $\mathbb{Z}_3 \subset SU(3)$ finding: the
conjecture's $SU(2)$ factor, like its $SU(3)$ factor, cannot be a
direct continuation of the framework's existing structure.  An
independent per-site $\mathbb{C}^2$ internal index --- naturally
identified with weak isospin --- is required for $SU(2)_W$ to
contribute a separate factor.  Without that extension the
framework's verified per-site carrier supports only $SO(3,1)
\times U(1)$ at dim 7, half the conjecture's dim 18.

\subsection{Direct-product structure on the extended $\mathbb{C}^{12}$: dim 18}
\label{subsec:direct_product_C12}

The previous two subsections identified the minimal per-site
extension required by the conjecture: the per-site amplitude
\begin{equation}
\label{eq:c12_amplitude}
\psi[\text{chir}, \text{iso}, \text{col}] \in
\mathbb{C}^2_\text{chir} \otimes \mathbb{C}^2_\text{iso} \otimes
\mathbb{C}^3_\text{col} \;\cong\; \mathbb{C}^{12},
\end{equation}
with the four conjecture factors acting on disjoint tensor slots:
\begin{align*}
SO(3,1):\;\;& J_a \otimes I_2 \otimes I_3,\quad
              K_a \otimes I_2 \otimes I_3,\\
SU(2)_W:\;\; & I_2 \otimes \tfrac{1}{2}\sigma_a \otimes I_3,\\
SU(3):\;\;   & I_2 \otimes I_2 \otimes \tfrac{1}{2}\lambda_a,\\
U(1):\;\;    & i\,I_{12}.
\end{align*}
The script \nolinkurl{automorphism_direct_product_extended.py}
verifies the within-factor commutation relations
($\mathfrak{so}(3,1)$, $\mathfrak{su}(2)$, $\mathfrak{su}(3)$ all
close correctly) and the pairwise commutation
$[\mathfrak{so}(3,1), \mathfrak{su}(2)] = [\mathfrak{so}(3,1),
\mathfrak{su}(3)] = [\mathfrak{su}(2), \mathfrak{su}(3)] = 0$ by
the tensor identity $[A \otimes I, I \otimes B] = 0$.  The
$\mathfrak{u}(1)$ generator $i I_{12}$ is central and commutes with
everything.

This establishes the direct-sum Lie algebra
\begin{equation}
\label{eq:direct_sum_18}
\mathfrak{so}(3,1) \oplus \mathfrak{su}(3) \oplus
\mathfrak{su}(2) \oplus \mathfrak{u}(1)
\;\subseteq\; \mathfrak{aut}(\mathcal{T}_\diamond^3,
\mathcal{A}=1)
\end{equation}
on the extended $\mathbb{C}^{12}$ carrier, with real dimension
$6 + 8 + 3 + 1 = 18$.  Both the dimension and the direct-sum
structure of the right-hand side of
Eq.~\eqref{eq:automorphism_conjecture} are now realised on the
extended amplitude.  What remains for the conjecture's
\emph{equality} is the question of whether the inclusion in
\eqref{eq:direct_sum_18} is proper --- whether the extended
$\mathrm{Aut}_\text{ext}$ admits any generator beyond the
direct-sum 18-dimensional Lie algebra constructed here.  That is
the open subproblem of
\autoref{subsec:open_subproblems}, item (i).

\subsection{Tick-rule consistency on the extended $\mathbb{C}^{12}$}
\label{subsec:tick_rule_consistency}

The bipartite chirality tick operator
\[
T_\text{chir} = \begin{pmatrix}
  i\sin(\delta\phi/2) & \cos(\delta\phi/2) \\
  \cos(\delta\phi/2)  & i\sin(\delta\phi/2)
\end{pmatrix}
\]
on $\mathbb{C}^2 = (\psi_R, \psi_L)$ is unitary and preserves
$\mathcal{A} = 1$ per Paper~I.  The natural extension to the
$\mathbb{C}^{12}$ carrier of
Eq.~\eqref{eq:c12_amplitude} is the trivial tensor extension
\begin{equation}
\label{eq:trivial_tensor_extension}
T_\text{ext} \;=\; T_\text{chir} \otimes I_2 \otimes I_3,
\end{equation}
acting as the existing tick rule on chirality and as the identity
on isospin and colour.  The script
\nolinkurl{tick_rule_extended_consistency.py} verifies four
structural properties:

\begin{enumerate}
\item \textbf{Unitarity}: $T_\text{ext}^\dagger T_\text{ext} = I_{12}$,
so $\mathcal{A}=1$ on $\mathbb{C}^{12}$ follows from $\mathcal{A}=1$
on the chirality $\mathbb{C}^2$ by linearity.
\item \textbf{Global $SU(2)_W$ commutativity}: $[T_\text{ext}, I_2
\otimes \tfrac{1}{2}\sigma_a \otimes I_3] = 0$ for $a = 1, 2, 3$.
\item \textbf{Global $SU(3)$ commutativity}: $[T_\text{ext}, I_2
\otimes I_2 \otimes \tfrac{1}{2}\lambda_a] = 0$ for $a = 1, \ldots, 8$.
\item \textbf{Bipartite parity}: with parity operator $P_\text{ext}
= \sigma_x \otimes I_2 \otimes I_3$, the identity $P_\text{ext}
T_\text{ext} P_\text{ext} = T_\text{ext}$ holds --- the chirality
parity that exchanges $\psi_R \leftrightarrow \psi_L$ is a
symmetry of the extended tick rule, and the isospin and colour
factors are parity-singlets in the trivial extension.
\end{enumerate}

The structural reading: the trivial tensor extension does no harm
to the framework's existing invariants.  $\mathcal{A}=1$, bipartite
parity, and the global $SU(2)_W \times SU(3)$ symmetry all
survive.  What this does \emph{not} establish is that the trivial
extension has physical content beyond six independent copies of the
existing chirality dynamics.  Promoting the global $SU(2)_W \times
SU(3)$ symmetry to a local gauge symmetry --- and identifying the
chirality-asymmetric coupling that the Standard Model demands --- is
the substance of the open subproblems in
\autoref{subsec:open_subproblems}, items (ii)--(iii).  The
established content here is precisely that the global symmetry is
\emph{available}; the gauge dynamics that turns it into physics is
the open work.

\subsection{Wilson plaquette gauge invariance on the bipartite lattice}
\label{subsec:wilson_plaquette}

Gauging the global $SU(2)_W \times SU(3)$ of
\autoref{subsec:tick_rule_consistency} introduces link variables
$U_i(x) \in G = U(1) \times SU(2)_W \times SU(3)$ on each
basis-vector edge, transforming as
\[
\psi(x) \to V(x)\psi(x),\qquad
U_i(x) \to V(x)\, U_i(x)\, V^\dagger(x + \mathbf{V}_i)
\]
under a local gauge transformation $V(x) \in G$.  The link-dressed
matter bilinear
\[
\psi^\dagger(x)\, U_i(x)\, \psi(x + \mathbf{V}_i)
\]
is gauge-invariant by direct substitution; the script
\nolinkurl{tick_rule_gauge_invariance.py} verifies this explicitly
for $U(1)$ and by sample check for $SU(2)$ link variables, with the
general non-abelian case (including $SU(3)$) following from
$V V^\dagger = I$.

The smallest closed gauge-invariant loop is structurally distinctive
on the bipartite octahedral lattice.  The three RGB basis vectors
satisfy $\mathbf{V}_1 + \mathbf{V}_2 + \mathbf{V}_3 = (1, 1, -1)
\neq 0$, so there are no triangular plaquettes; the smallest
closed loop respecting the RGB/CMY alternation is the four-link
``bipartite plaquette''
\[
x \to x + \mathbf{V}_1 \to x + \mathbf{V}_1 - \mathbf{V}_2 \to
x - \mathbf{V}_2 \to x,
\]
traversing the basis vectors in the cyclic order $\mathbf{V}_1,
-\mathbf{V}_2, -\mathbf{V}_1, \mathbf{V}_2$.  The plaquette variable
$W_p = U_1(x)\, U_{-2}(x + \mathbf{V}_1)\, U_{-1}(x + \mathbf{V}_1
- \mathbf{V}_2)\, U_2(x - \mathbf{V}_2)$ has gauge-invariant trace
by the cyclic property of $\mathrm{Tr}$ --- the gauge
transformations $V(x), V(x + \mathbf{V}_1), V(x + \mathbf{V}_1 -
\mathbf{V}_2), V(x - \mathbf{V}_2)$ at the four vertices cycle and
cancel.  This holds for any $SU(N)$ gauge group, in particular for
the $SU(2)_W$ and $SU(3)$ of the conjecture.

The Wilson action template
\begin{equation}
\label{eq:wilson_action_template}
S_W \;=\; \frac{1}{g^2} \sum_p \left(1 - \frac{1}{N}
\mathrm{Re}\,\mathrm{Tr}\, W_p\right)
\end{equation}
is therefore gauge-invariant on the bipartite octahedral lattice
by inspection.  The framework's existing $U(1)$ induced-action
calculation (Paper~I, Appendix~B) is the worked example of how the
bare coupling $1/g^2$ is extracted at the lattice scale; the
template \eqref{eq:wilson_action_template} extends to $SU(2)_W$ and
$SU(3)$ with matrix-valued link variables and no other structural
change.  Computing the explicit $1/g^2(a)$ prefactor for the
non-abelian cases is the substance of the open subproblem of
\autoref{subsec:open_subproblems}, item (iii).

\subsection{Necessity vs sufficiency: established containment $\supseteq$, conjectured equality $=$}
\label{subsec:established_necessity_vs_sufficiency}

The six PASS sub-claims of
Eq.~\eqref{eq:automorphism_conjecture} combine as follows; the
seventh PASS row (SU(3) representation branch consistency) narrows
the colour interpretation that the combination assumes (see
\nolinkurl{notes/su3_branch_consistency.md}).
\autoref{subsec:discrete_spatial} establishes the discrete spatial
backbone $O_h$.  \autoref{subsec:rgb_z3_obstruction} and
\autoref{subsec:lorentz_u1_on_C2} together force the per-site
internal extension $\mathbb{C}^2_\text{iso} \otimes
\mathbb{C}^3_\text{col}$: $SU(2)_W$ cannot share the chirality
$\mathbb{C}^2$ with $SO(3,1)$ rotations, and $SU(3)$ cannot be a
Lie-theoretic continuation of the discrete RGB symmetry.
\autoref{subsec:direct_product_C12} verifies that the extended
$\mathbb{C}^{12}$ carrier supports the four conjecture factors as
a direct-sum 18-dimensional Lie algebra with the structure
constants of $\mathfrak{so}(3,1) \oplus \mathfrak{su}(3) \oplus
\mathfrak{su}(2) \oplus \mathfrak{u}(1)$.
\autoref{subsec:tick_rule_consistency} verifies that the bipartite
tick rule extends consistently to $\mathbb{C}^{12}$ as a global
symmetry of the dynamics.  \autoref{subsec:wilson_plaquette}
verifies that the global symmetry can be promoted to a local gauge
symmetry with a gauge-invariant Wilson plaquette on the smallest
non-trivial closed loop the bipartite lattice admits; the explicit
$1/g^2$ prefactor that turns the gauge-invariant plaquette into a
full Wilson action remains open (item (iii) below).

The conjunction is the inclusion of
Eq.~\eqref{eq:direct_sum_18}: the conjecture's right-hand side is
contained in the lattice's automorphism algebra,
\begin{equation}
\label{eq:established_containment}
\mathrm{Aut}(\mathcal{T}_\diamond^3, \mathcal{A}=1)
\;\supseteq\;
SO(3,1) \times SU(3) \times SU(2) \times U(1),
\end{equation}
with both sides finite-dimensional Lie groups of real dimension
$\geq 18$ on the extended $\mathbb{C}^{12}$ carrier.

The containment in \eqref{eq:established_containment} is necessary
but not sufficient for the conjectured equality
\eqref{eq:automorphism_conjecture}.  Three structural questions
remain open, in the order they enter the equality:
\begin{enumerate}
\item \emph{Equality vs proper containment in
$\mathrm{Aut}_\text{ext}$}: whether the extended automorphism
algebra admits any generator beyond the 18 already constructed.
This is a finite Lie-algebra enumeration question constrained by
commutation with the established factors.

\item \emph{Chirality coupling alignment} (partially resolved,
characterisation outcome): whether the framework's bipartite parity
(the RGB/CMY $\mathbb{Z}_2$) is the SM's chirality projector, and
whether the proposed $SU(2)_W$ on the isospin $\mathbb{C}^2$
couples to $\psi_R$ and $\psi_L$ asymmetrically as the Standard
Model demands.  Under Branch~A SU(3) and the existing tick rule,
the unique linear bipartite parity $P = \sigma_x \otimes I_2 \otimes
I_3$ anticommutes with $\gamma_5 = \sigma_z \otimes I_2 \otimes I_3$
and swaps $\psi_L \leftrightarrow \psi_R$ --- it IS spatial parity,
not chirality projection.  The proposed $SU(2)_W$ commutes with $P$
(vector-like, not the SM's chiral coupling).  The lattice's
bipartite $\mathbb{Z}_2$ and the SM's chirality $\mathbb{Z}_2$ are
distinct orthogonal Bloch involutions on the chirality
$\mathbb{C}^2$.  The trivial tensor extension
\eqref{eq:trivial_tensor_extension} therefore predicts a precise
\emph{characterisation} of the parity / chirality relationship
rather than an alignment or a vacuous obstruction.  The follow-on
question --- whether a modified antilinear tick rule of the form
$T_\text{ext}^B = is \cdot I_{12} + c \cdot (\sigma_x \otimes I_2
\otimes C) \circ K$ can admit the Branch~B SU(3) interpretation and
deliver SM-style $CP$ --- is settled negatively: no $3 \times 3$
matrix $C$ anticommutes with every Gell-Mann, by direct argument on
$\lambda_3$ and $\lambda_8$, equivalently because $\mathbf{3}$ and
$\bar{\mathbf{3}}$ are inequivalent irreps of SU(3) and the unique
antilinear intertwiner between them (global colour conjugation)
fails to satisfy the chirality-dependent commutativity Branch~B
demands.  The lattice's bipartite parity is therefore intrinsically
\emph{linear} (spatial parity, no charge conjugation), and the
framework cannot incorporate SM-style $CP$ as a discrete symmetry
of any natural tick modification.

\item \emph{Explicit $1/g^2$ prefactors for the non-abelian Wilson
actions}: generalising Paper~I's $U(1)$ induced-gauge-action
calculation to $SU(2)_W$ and $SU(3)$ with matrix-valued links, and
checking the resulting bare couplings against the measured
Standard-Model values at the lattice scale.
\end{enumerate}
These three questions are the substance of
\autoref{subsec:open_subproblems}.  Their joint resolution decides
between the conjecture's two terminal states: equality (Paper~II
publishes the derivation of the SM gauge group) or precise
characterisation of the obstruction (Paper~II publishes the
framework's verified scope plus the structural reason the equality
fails).  Both are concrete results; the established containment
\eqref{eq:established_containment} is what makes either of them
publishable.
